\cotefile{15}
\scanpage{87}
\margin{VI. 5}

$\rho_{\widetilde A}:\widetilde A\longrightarrow j_*j^*(\widetilde A)$ is an isomorphism, i.e. $j$ satisfies condition (b) of Section V as well as condition (a).

An important property of the integrally closed noetherian domain $A$ which will be invoked in several succeeding proofs is the following theorem: if $M$ and $N$ are finitely generated $A$-modules and there is an isomorphism $f:M\longrightarrow N$ such that $f_{\mathfrak p}:M_{\mathfrak p}\longrightarrow N_{\mathfrak p}$ is an isomorphism for every $\mathfrak p$ in $P$, then $f^{**}:M^{**}\longrightarrow N^{**}$ is an isomorphism (where $M^{**}$ is $\operatorname{Hom}_A(\operatorname{Hom}_A(M,A),A)$, the double dual of $M$). A proof of this theorem can be found in (amongst other sources) the thesis of the author, Chapter II, Lemma 0.1.

The following theorems will enable us to reinterpret the groups in the exact sequence of Theorem V.12 in terms of algebras and modules over $A$.

THEOREM VI.5. Let $Y$ be a subset of $X$ containing $P$. Then if $\mathcal M$ is locally free of finite rank over $(Y,j^*(\widetilde A))$, there is a finitely generated reflexive module $N$ such that $j^*(\widetilde N)\approx\mathcal M$.

\uline{Proof:} We know that the finitely generated module $N=\Gamma(U,j_*(\mathcal M)|U)=\Gamma(Y,\mathcal M)$ has the property that $j^*(\widetilde N)\approx\mathcal M$, where $U$, as defined in the proof of Theorem VI.1, is an open set containing $Y$. We will show that $N$ is reflexive. Let $N^{**}=\operatorname{Hom}_A(\operatorname{Hom}_A(N,A),A)$. Then there is a homomorphism $\delta:N\longrightarrow N^{**}$, giving rise to the sheaf

\mspage{86}
